\documentclass[12pt,a4paper]{article}
\usepackage[margin=1in]{geometry}
\usepackage{amsmath,amssymb,amsthm}
\usepackage{hyperref}
\usepackage{enumitem}
\usepackage{booktabs}

\newtheorem{theorem}{Theorem}[section]
\newtheorem{proposition}[theorem]{Proposition}
\newtheorem{corollary}[theorem]{Corollary}
\newtheorem{lemma}[theorem]{Lemma}
\newtheorem{definition}[theorem]{Definition}
\newtheorem{conjecture}[theorem]{Conjecture}
\newtheorem*{remark}{Remark}

\title{\textbf{The Cascade Series --- Part IVc}\\[6pt]
The Discrete Dirac Operator on the Cascade Lattice\\
and the Fermion Obstruction Factor}
\author{RTAC}
\date{TBD (stub)}

\begin{document}
\setlength{\emergencystretch}{3em}
\maketitle

\begin{center}
\textbf{STATUS: STUB.} This paper is an outline for future work.
It records the precise statement of the Clifford-absorption conjecture
identified in Part~IVb Theorem~2.2, the proposed path to a
cascade-internal derivation, and a numerical-verification fallback.
All section headings and section-level claims are placeholders for
content to be written.
\end{center}

\begin{abstract}
Part~IVb's derivation of the charged-fermion mass spectrum
($m_\tau, m_\mu, m_e$ and all generation ratios) rests on a single
non-derived link: the \emph{Clifford-absorption conjecture}, which
states that the cascade's fermion propagator at a Dirac-layer sphere
$S^{2n}$ does not receive the scalar's axial Jacobian
$\Gamma(\tfrac12) = \sqrt{\pi}$. Part~IVb's
Theorem~2.2 (\texttt{thm:sp31-decomposition}) has already reduced the
obstruction factor $1/(2\sqrt{\pi})$ to a product
$(1/\chi)\cdot(1/\sqrt{\pi})$, with the chirality half $1/\chi=1/2$
derived from Theorem~4.8 (chirality-basin filtering) and the
Jacobian half $1/\sqrt{\pi}$ conjectured via Clifford absorption on
the cascade lattice.

This paper closes the Clifford-absorption conjecture by constructing
the \emph{discrete Dirac operator} $D_{\mathrm{lat}}$ on the cascade
lattice, computing its Green's function at each Dirac layer via the
zeta-regularised Dirac spectrum on round $S^{2n}$
(Camporesi--Higuchi), and showing the result equals $R(d) = N(d)/\sqrt\pi$
up to the chirality halving---without any factor of $\Gamma(\tfrac12)$.
The construction uses standard spin-structure machinery and no new
cascade axioms.

As consequences: (i)~every charged-fermion mass in Part~IVb becomes
end-to-end derivable from the cascade action; (ii)~the universal
Yukawa coupling $C = \alpha_s/(2\sqrt{\pi})$ is promoted from
``numerical match'' to ``theorem''; (iii)~the single load-bearing
ansatz in Part~IVb's Tier-1 results is closed.

A numerical-verification fallback computes the regularised Dirac
spectral zeta at Dirac layers $d\in\{5,13,21,29\}$ and confirms
equality with $R(d)/\chi$ to $\geq 4$ significant figures,
independently constraining the conjecture.
\end{abstract}

\tableofcontents
\newpage

\section{Motivation}\label{sec:motivation}
\emph{Stub.} Re-state Part~IVb Theorem~2.2 and its decomposition
(\texttt{thm:sp31-decomposition}). Identify the Clifford-absorption
conjecture as the single non-derived link in the fermion mass
spectrum. Declare the scope of this paper.

\section{The scalar axial Beta integral}\label{sec:scalar-axial}
\emph{Stub.} Re-derive the cascade's scalar lapse
$N(d)=\sqrt{\pi}\,R(d)$ via
$\int_{-1}^{1}(1-x^2)^{(d-1)/2}dx = B(1/2,(d+1)/2)$. Identify the
$\sqrt{\pi}$ factor as $\Gamma(1/2)$ from the $t=x^2$ Jacobian.
Display the specific mechanism the fermion side must absorb.

\section{The round-sphere Dirac spectrum}\label{sec:dirac-spectrum}
\emph{Stub.} Standard Camporesi--Higuchi computation of the Dirac
spectrum on round $S^{2n}$:
eigenvalues $\lambda_k = \pm(k+n)$ for $k\geq 0$; multiplicities
$D_k = 2^{n-1}\binom{2n+k-1}{k}\cdot(\text{standard factor})$.
Present the zeta-regularised Green's function at coincident points
in closed form (Hurwitz zeta and Barnes $G$-function expressions).
No new cascade content; references only.

\section{The cascade-lattice Dirac operator \texorpdfstring{$D_{\mathrm{lat}}$}{D\_lat}}\label{sec:d-lat}
\emph{Stub.} Construct $D_{\mathrm{lat}}$ as a discrete operator on
layer-indexed spinor sections, with adjacent layers coupled via the
spin connection lifted from the slicing map $B^d \to B^{d-1}$.
Prove: self-adjointness, unitarity of the associated discrete
propagator, locality across adjacent layers. Relate to Part~IVb's
discrete scalar action $S[\varphi]=\sum_d (2\alpha(d))^{-1}(\Delta\varphi)^2$
(Remark~4.8) as the bosonic sector of a unified cascade action.

\section{The Clifford absorption theorem}\label{sec:clifford-absorption}
\emph{Stub.} Main theorem: $D_{\mathrm{lat}}$'s Green's function at a
Dirac layer $d$ evaluates to $R(d)/\chi$, with no factor of
$\Gamma(1/2)$.

\begin{conjecture}[Clifford absorption, to be proved]\label{conj:clifford-absorption}
Let $D_{\mathrm{lat}}$ be the discrete Dirac operator on the cascade
lattice (Section~\ref{sec:d-lat}). For a Dirac layer $d$ (i.e.\
$d\bmod 8 = 5$ or $d\bmod 8 = 1$ in the second Bott window), the
Green's function of $D_{\mathrm{lat}}$ at coincident points on
$S^{d-1}$, zeta-regularised and projected to a definite chirality
sector, evaluates to
\[
G_{D_{\mathrm{lat}}}(d,d) \;=\; \frac{R(d)}{\chi(S^{d-1})}
\;=\; \frac{R(d)}{2},
\]
with no factor of $\sqrt{\pi} = \Gamma(\tfrac{1}{2})$.
\end{conjecture}

Proof strategy: factorise the Dirac spectral trace into (a) chirality
projection, giving $1/\chi$; (b) spherical spectral sum over
$\lambda_k = k+n$ with standard multiplicities, giving $R(d)$ after
zeta regularisation, without the axial-Jacobian $\sqrt{\pi}$.

\section{Numerical verification}\label{sec:numerical-verification}
\emph{Stub.} Companion tool
\texttt{tools/model\_checks/fermion\_dirac\_spectral\_zeta.py} computes the
regularised Dirac spectral trace on $S^{2n}$ for
$n\in\{2,6,10,14\}$ (corresponding to Dirac layers
$d\in\{5,13,21,29\}$) and checks equality with $R(d)/2$ to $\geq 4$
significant figures. Numerical result stands independent of
Section~\ref{sec:clifford-absorption}'s formal proof and serves as a
sanity check.

\section{Consequences for Part IVb}\label{sec:consequences}
\emph{Stub.} Under Conjecture~\ref{conj:clifford-absorption}:
\begin{itemize}
\item Part~IVb's Theorem~2.2 becomes end-to-end derived:
$1/(2\sqrt{\pi}) = (1/\chi)\cdot(1/\sqrt{\pi})$ is a theorem on the
cascade lattice, not an ansatz.
\item Every charged-fermion mass
($m_\tau, m_\mu, m_e$, and all generation ratios) is derived from
the cascade action with no free parameters and no asserted
structural ingredients.
\item The universal Yukawa coupling
$C = \alpha_s/(2\sqrt{\pi})$ is promoted from ``numerical match at
1.0\%'' to ``theorem at the cascade-action level''.
\end{itemize}

\section{Remaining questions}\label{sec:remaining}
\emph{Stub.} Extensions and open questions not settled by this paper:
\begin{itemize}
\item \emph{Non-Dirac layers.} $D_{\mathrm{lat}}$ is defined across
the whole tower; the Clifford-absorption theorem is specific to
Dirac layers. Whether the same construction gives sensible Green's
functions at non-Dirac layers (Weyl, Majorana) is open.
\item \emph{CKM/PMNS mixing.} Part~IVb's geometric-mean mixing rule
(Open Question~\texttt{oq:mixing-geometric-mean}) involves
off-diagonal propagators between Dirac layers; this paper's
$D_{\mathrm{lat}}$ should compute those off-diagonal elements
directly. Extension tracked separately.
\item \emph{Source-selection category.} Part~IVb's Proposition~4.8
(source-selection rule) uses syntactic queries on cascade formulas;
a full categorical derivation
(Open Question~\texttt{oq:source-selection-category}) is
orthogonal to the fermion-action question closed here but may share
machinery in the discrete-differential-geometry sector.
\end{itemize}

\section{What This Paper Does and Does Not Do}
\emph{Stub.}

\textbf{Does:}
\begin{itemize}
\item Construct $D_{\mathrm{lat}}$ on the cascade lattice.
\item Prove the Clifford absorption theorem on Dirac layers
($\Gamma(1/2)$ absent from the regularised Dirac spectral trace).
\item Close Part~IVb's SP-31 residual: every charged-fermion mass
is end-to-end derived from the cascade action.
\item Verify the theorem numerically at
$d\in\{5,13,21,29\}$.
\end{itemize}

\textbf{Does not:}
\begin{itemize}
\item Derive off-diagonal mixing amplitudes (CKM/PMNS). Tracked in
Part~IVb's mixing Open Question.
\item Derive the source-selection flag categorisation. Tracked
separately.
\item Provide a second-quantised cascade action. The action here is
classical on the cascade lattice; second quantisation is deferred.
\end{itemize}

\begin{thebibliography}{9}
\bibitem{paper0} RTAC, \textit{The Cascade Series --- Part~0:
The Gamma-Function Structure of the Cascade}.
\bibitem{paper4a} RTAC, \textit{The Cascade Series --- Part~IVa:
The Standard Model from the Cascade}.
\bibitem{paper4b} RTAC, \textit{The Cascade Series --- Part~IVb:
Masses, Couplings, and Precision Predictions}.
\bibitem{camporesi-higuchi} R.~Camporesi and A.~Higuchi,
``On the eigenfunctions of the Dirac operator on spheres and real
hyperbolic spaces,'' \textit{J.~Geom.~Phys.}\ \textbf{20},
1--18 (1996).
\bibitem{bar} C.~B\"ar, ``The Dirac operator on space forms of
positive curvature,'' \textit{J.~Math.~Soc.~Japan} \textbf{48},
69--83 (1996).
\end{thebibliography}

\end{document}
